\chapter{Proposal Action Register}
\label{app:action_reg}

This appendix reproduces the committee's written feedback from the FYP Proposal Defense I for Project \textbf{F26-012-FORGE} in full, followed by the resolution register for Iterations 1 and 2. The register groups the feedback into ten directives; the requests to explore recent RL techniques and to evaluate engagement and satisfaction are covered under rows 6 and 5.

\section*{Official Evaluation Committee Review}
\begin{itemize}
    \item \textbf{Project Code:} F26-012-FORGE
    \item \textbf{Decision:} Approved with Minor Modifications
    \item \textbf{Issuing Authority:} FYP Committee, School of Computing, FAST-NUCES, Islamabad
\end{itemize}

\subsection*{Comments / Recommended Modifications (verbatim)}
\begin{quote}\small
Good project. Good presentation and demonstration. Demonstrates satisfactory understanding of the project concepts and approach. The game's aesthetics should be improved. They should study related literature review deeply as the game is based on research. They only covered one aspect. Students need to explore recent RL techniques, as well as other areas of the latest research. Students also need to work on the aesthetics of the application, as in the game industry there is a lot of competition with respect to features as well as aesthetics. Overall concept is fine. Focus also on aesthetics; also keep a variety of attacks for enemy and player powers. Your game should also give a good, playable, engaging experience without AI. Write the scope clearly regarding what you will do and what you will not. The final game/product should have a clear, measurable way of evaluating your RL model and game. Extend the literature review to comprehensively cover Dynamic Difficulty Adjustment (DDA), in addition to procedural content generation. The review should identify established player-performance, engagement, satisfaction, and difficulty measures and use these to refine the proposed player model and evaluation strategy. Strengthen the technical justification for PPO by comparing it with appropriate alternative reinforcement-learning approaches and explaining why PPO is suitable for the expected gradual and abrupt changes in player performance. Similarly, define and validate the judge model using empirical player data rather than primarily manually assigned difficulty weights. Clearly distinguish procedural level generation from difficulty adaptation and demonstrate the benefit of the closed-loop adaptive approach over a strong non-adaptive procedural-generation baseline. Evaluation should consider not only difficulty matching but also player engagement and satisfaction. Further specify measurable acceptance criteria, including convergence time, solvability, difficulty-prediction error, layout diversity, adaptation stability, and player-experience measures. Also document risks, fallback strategies, and ethical considerations related to player profiling and telemetry data.
\end{quote}

\section*{Detailed Resolution Register}
\begin{longtable}{p{0.05\textwidth}p{0.26\textwidth}p{0.32\textwidth}p{0.12\textwidth}p{0.19\textwidth}}
\toprule
\textbf{No.} & \textbf{Committee Directive} & \textbf{Action Taken in Code \& Architecture} & \textbf{Status} & \textbf{Evidence} \\
\midrule
1 & ``The game's aesthetics should be improved. They should also work on the aesthetics of the application, as in the game industry there is a lot of competition with respect to features as well as aesthetics.'' & Replaced geometric circles with 2D pixel-art top-down sprites; integrated trauma-based camera shake (decay rate $3.6$\,s$^{-1}$), animated floating combat text (\texttt{DamageFloater.tscn}), and Cinzel/Inter typography. & Completed & §\ref{sec:iter1_features}, Figs.~\ref{fig:app_sprites}, \texttt{Main.gd} \\
\midrule
2 & ``Focus also on aesthetics; also keep a variety of attacks for enemy and player powers.'' & Engineered Ranged Pyromancer with kiting AI and telegraphed fireball projectiles; added player Spacebar dodge roll dash with invulnerability frames and cooldown. & Completed & \texttt{Enemy.gd}, \texttt{Player.gd}, \texttt{ProjectileFireball.gd} \\
\midrule
3 & ``Your game should also give a good, playable, engaging experience without AI.'' & Developed an in-engine toggle in the Defense Telemetry Overlay enabling real-time switching to a deterministic $+8\%$ per-floor stat progression static baseline. & Completed & §\ref{sec:iter1_features}, \texttt{DefenseOverlay.gd} \\
\midrule
4 & ``Write the scope clearly regarding what you will do and what you will not.'' & Formulated a formal three-tier Scope Boundary Matrix defining strict Mid-1 deliverables, Final-1/FYP-2 commitments, and explicitly out-of-scope boundaries. & Completed & Table~\ref{tab:scope_matrix}, §\ref{sec:scope} \\
\midrule
5 & ``The final game/product should have a clear, measurable way of evaluating your RL model and game.'' & Formulated 8 quantitative evaluation metrics (solvability, latency, DER, $CR_0$ convergence) and standardized psychometric instruments (GEQ, PENS). Convergence time, layout diversity and adaptation stability are not yet defined. & Partially completed & Table~\ref{tab:success_criteria}, Chapter~\ref{chap:testing} \\
\midrule
6 & ``Extend the literature review to comprehensively cover Dynamic Difficulty Adjustment (DDA)... identify established player-performance, engagement, satisfaction, and difficulty measures.'' & Conducted a structured review of 40 sources (IEEE ToG, IEEE CoG, ACM CHI, Springer MTAP and others), including recent RL-based and player-adaptive work; grouped DDA into 3 paradigms; synthesized flow and engagement measures. & Completed & Chapter~\ref{chap:research}, Table~\ref{tab:lit_comparison} \\
\midrule
7 & ``Strengthen the technical justification for PPO by comparing it with appropriate alternative reinforcement-learning approaches and explaining why PPO is suitable for gradual and abrupt changes.'' & Evaluated PPO against SAC, DQN, and Contextual Bandits; explained how PPO's clipped surrogate objective keeps training updates small under abrupt performance changes, and its limits (no strict ratio bound or monotonic guarantee); deployment-time smoothness is measured (§\ref{sec:policy_eval}). & Completed & Section~\ref{sec:ppo_proof}, Table~\ref{tab:rl_comparison} \\
\midrule
8 & ``Define and validate the judge model using empirical player data rather than primarily manually assigned difficulty weights.'' & Hand-weighted linear proxy (not yet validated) active for Mid-1; supervised Gradient Boosted Regressor with TreeSHAP and 25-player empirical corpus formally scheduled for FYP-1 Final (Gate 3 Exit, Weeks 11--14). & Scheduled for Final-1 & Section~\ref{sec:domain}, Chapter~\ref{chap:fyp2} \\
\midrule
9 & ``Clearly distinguish procedural level generation from difficulty adaptation and demonstrate the benefit of the closed-loop adaptive approach over a strong non-adaptive procedural baseline.'' & Decoupled architecture into Stage 1 (connected spatial topology, BFS-checked in CI) and Stage 2 (closed-loop PPO combat adapter); provided the static roguelike baseline. The comparison demonstrating the benefit over it is scheduled for FYP-1 Final. & Partially completed & Section~\ref{sec:proposed}, Chapter~\ref{chap:design} \\
\midrule
10 & ``Document risks, fallback strategies, and ethical considerations related to player profiling and telemetry data.'' & Developed zero-crash heuristic fallback in \texttt{PPOPolicy.gd}; implemented client-side zero-PII local data governance; formulated risk register. & Completed & Section~\ref{sec:ethics}, Section~\ref{sec:risk} \\
\bottomrule
\end{longtable}
